\subsection{典范内射消解}

设 F 为 A-模 \(\mathrm{Hom}_{\mathbf{Z}}(\mathbf{A},\mathbf{Q} / \mathbf{Z})\) ；对于每个 A-模 M，我们设 \(\mathrm{I}^0 (\mathbf{M}) = \mathrm{F}^{\mathrm{Hom}_{\mathbf{A}}(\mathbf{M},\mathbf{F})}\) 并用 \(e_{\mathbf{M}}:\mathbf{M}\to \mathrm{I}^{0}(\mathbf{M})\) 表示将 \(m\in \mathbf{M}\) 映到族 \((\varphi (m))_{\varphi \in \mathrm{Hom}_{\mathbf{A}}(\mathbf{M},\mathbf{F})}\) 的同态。根据X，第19页，推论2， \(\mathrm{I}^0 (\mathbf{M})\) 是一个内射 A-模，且 \(e_{\mathbf{M}}\) 是内射的。设 \(\mathbf{K}^0 (\mathbf{M}) = \mathbf{Coker}\) \(e_{\mathbf{M}}\) 并记 \(q_{\mathbf{M}}:\mathrm{I}^0 (\mathbf{M})\to \mathrm{K}^0 (\mathbf{M})\) 为典范投影。因此我们有一个正合序列

\[
0 \longrightarrow \mathrm{M} \stackrel {e _ {\mathrm{M}}} {\longrightarrow} \mathrm{I} ^ {0} (\mathrm{M}) \stackrel {q _ {\mathrm{M}}} {\longrightarrow} \mathrm{K} ^ {0} (\mathrm{M}) \longrightarrow 0.
\]

我们定义一个分次A-模I(M)，设 \(\mathbf{I}^n (\mathbf{M}) = 0\) 于 \(n < 0\)，并对整数 \(n > 0\) 归纳地设

\[
\mathrm{I} ^ {n} (\mathbf {M}) = \mathrm{I} ^ {0} \left(\mathrm{K} ^ {n - 1} (\mathbf {M})\right), \quad \mathrm{K} ^ {n} (\mathbf {M}) = \mathrm{K} ^ {0} \left(\mathrm{K} ^ {n + 1} (\mathbf {M})\right).\tag{9}
\]

我们由下式定义 A-同态 \(\delta_{\mathbf{M}}^{n}:\mathrm{I}^{n}(\mathbf{M})\to \mathrm{I}^{n + 1}(\mathbf{M})\)

\[
\left\{ \begin{array}{l l} \delta_ {\mathbf {M}} ^ {n} = 0, & n <   0, \\ \delta_ {\mathbf {M}} ^ {0} = e _ {\mathbf {K} ^ {0} (\mathbf {M})} \circ q _ {\mathbf {M}}, \\ \delta_ {\mathbf {M}} ^ {n} = e _ {\mathbf {K} ^ {n} (\mathbf {M})} \circ q _ {\mathbf {K} ^ {n - 1} (\mathbf {M})}, & n > 0. \end{array} \right.\tag{10}
\]

由构造，我们有一个正合序列

\[
0 \longrightarrow \mathrm{M} \xrightarrow {e _ {\mathrm{M}}} \mathrm{I} ^ {0} (\mathrm{M}) \xrightarrow {\delta_ {\mathrm{M}} ^ {0}} \dots \longrightarrow \mathrm{I} ^ {n} (\mathrm{M}) \xrightarrow {\delta_ {\mathrm{M}} ^ {n}} \mathrm{I} ^ {n + 1} (\mathrm{M}) \longrightarrow \dots ,
\]

因此，如果把 \(e_{M}\) 延拓为复形的态射

\[
e _ {\mathbf {M}}: \mathbf {M} \rightarrow (\mathrm{I} (\mathbf {M}), \delta_ {\mathbf {M}}),
\]

我们得到M的一个内射分解，称为M的典范内射分解。设 \(f: \mathbf{M} \to \mathbf{N}\) 是 A-模的同态。用 \(\mathrm{I}^{0}(f)\) 表示从 \(\mathrm{I}^{0}(\mathbf{M}) = \mathrm{F}^{\mathrm{Hom}_{\mathbf{A}}(\mathbf{M},\mathbf{F})}\) 到 \(\mathrm{I}^{0}(\mathbf{N}) = \mathrm{F}^{\mathrm{Hom}_{\mathbf{A}}(\mathbf{N},\mathbf{F})}\) 的、将族 \((x_{\varphi})_{\varphi \in \mathrm{Hom}_{\mathbf{A}}(\mathbf{M},\mathbf{F})}\) 映到族 \((x_{\psi \circ f})_{\psi \in \mathrm{Hom}_{\mathbf{A}}(\mathbf{N},\mathbf{F})}\) 的同态。我们有：

\[
\mathrm{I} ^ {0} (f) \circ e _ {\mathbf {M}} = e _ {\mathbf {N}} \circ f.\tag{11}
\]

因此， \(\mathrm{I}^0 (f)\) 诱导一个同态 \(\mathbf{K}^0 (f):\mathbf{K}^0 (\mathbf{M})\to \mathbf{K}_-^0 (\mathbf{N})\) 并且我们有

\[
\mathrm{K} ^ {0} (f) \circ q _ {\mathbf {M}} = q _ {\mathbf {N}} \circ \mathrm{K} ^ {0} (f).\tag{12}
\]

令 \(\mathrm{I}^n(f) = 0\) 当 \(n < 0\) 并通过对整数 \(n > 0\) 的归纳，定义同态 \(\mathrm{I}^n(f): \mathrm{I}^n(\mathbf{M}) \to \mathrm{I}^n(\mathbf{N})\) 与 \(\mathrm{K}^n(f): \mathrm{K}^n(\mathbf{M}) \to \mathrm{K}^n(\mathbf{N})\) 如下：

\[
\left\{ \begin{array}{l} \mathrm{I} ^ {n} (f) = \mathrm{I} ^ {0} (\mathrm{K} ^ {n - 1} (f)) \\ \mathrm{K} ^ {n} (f) = \mathrm{K} ^ {0} (\mathrm{K} ^ {n - 1} (f)). \end{array} \right.\tag{13}
\]

命题5. --- I(f) : I(M) → I(N) 是A-模复形的一个态射；我们有 I(f) ∘ e\_M = e\_N ∘ f。

这是以与命题4类似的方式证明的。

我们有

\[
\mathrm{I} (1 _ {\mathbf {M}}) = 1 _ {\mathrm{I} (\mathbf {M})}\tag{14}
\]

且对于每个A-模同态 \(g: N \rightarrow P\)

\[
\mathrm{I} (g \circ f) = \mathrm{I} (g) \circ \mathrm{I} (f).\tag{15}
\]

注记. --- 如果 \(f, g \in \operatorname{Hom}_{\mathbf{A}}(\mathbf{M}, \mathbf{N})\) ，我们没有 \(\mathrm{I}(f + g) = \mathrm{I}(f) + \mathrm{I}(g)\)。然而，这两个态射依X，第49页，命题3 bis是同伦的。

若M是一个右A-模，我们设 \(\mathrm{I}(\mathbf{M}) = \mathrm{I}(\mathbf{M}^{\circ})^{\circ}\) ；它被称为M的典范内射分解，并且我们有

\[
\mathrm{I} (\mathbf {M} ^ {\circ}) = \mathrm{I} (\mathbf {M}) ^ {\circ}.
\]

